\documentclass[../../main.tex]{subfiles}
\begin{document}

\begin{flushright}
  \footnotesize\textit{【自動文字起こし・要確認】}
\end{flushright}

{\bf [解答]}
$n$秒後に異なる点にいる確率を$Q(n)$とし，同じ点にいる確率を$P(n)$とする．

$2$つの粒子が同じ点にいる状態からは，各粒子が独立に残り$2$頂点のいずれかへ等確率で移動するので，
次に同じ点にいる確率は$\dfrac{1}{2}\cdot\dfrac{1}{2}+\dfrac{1}{2}\cdot\dfrac{1}{2}=\dfrac{1}{2}$である．

$2$つの粒子が異なる点にいる状態からは，同様に考えて，次に同じ点にいる確率は$\dfrac{1}{2}\cdot\dfrac{1}{2}=\dfrac{1}{4}$である．

\begin{figure}[htb]
  \centering
  \begin{tikzpicture}[>=stealth, every node/.style={font=\small}]
    \node[draw=blue!85!black, rounded corners=3pt, fill=blue!9,
    minimum width=2.4cm, minimum height=0.8cm] (S0) at (0,0) {同じ点 $P(n)$};
    \node[draw=teal!80!black, rounded corners=3pt, fill=teal!9,
    minimum width=2.4cm, minimum height=0.8cm] (D0) at (0,-2.2) {異なる点 $Q(n)$};
    \node[draw=blue!85!black, rounded corners=3pt, fill=blue!9,
    minimum width=2.6cm, minimum height=0.8cm] (S1) at (5.2,0) {同じ点 $P(n+1)$};
    \node[draw=teal!80!black, rounded corners=3pt, fill=teal!9,
    minimum width=2.6cm, minimum height=0.8cm] (D1) at (5.2,-2.2) {異なる点 $Q(n+1)$};
    \draw[->, line width=1.2pt, blue!85!black]
    (S0)--node[above] {$1/2$} (S1);
    \draw[->, line width=1.2pt, blue!85!black]
    (S0.south east)--node[pos=0.23, above right] {$1/2$} (D1.north west);
    \draw[->, line width=1.2pt, teal!80!black]
    (D0.north east)--node[pos=0.77, below right] {$1/4$} (S1.south west);
    \draw[->, line width=1.2pt, teal!80!black]
    (D0)--node[below] {$3/4$} (D1);
  \end{tikzpicture}
  \caption{2つの粒子が同じ点にいる状態と異なる点にいる状態の遷移}
\end{figure}

上図から
\begin{align}
  P(n+1)&=\dfrac{1}{2}P(n)+\dfrac{1}{4}Q(n)\\
  &=\dfrac{1}{2}P(n)+\dfrac{1}{4}(1-P(n))\\
  &=\dfrac{1}{4}P(n)+\dfrac{1}{4}
\end{align}

$P(1)=\dfrac{1}{2}$だから，くり返し用いて，
\begin{align}
  P(n)&=\left(\dfrac{1}{4}\right)^{n-1}\left(\dfrac{1}{2}-\dfrac{1}{3}\right)+\dfrac{1}{3}\\
  &=\dfrac{1}{3}\left\{1+\dfrac{1}{2}\left(\dfrac{1}{4}\right)^{n-1}\right\}．
\end{align}

したがって，$p(n)=\dfrac{1}{3}\left\{1+\dfrac{1}{2}\left(\dfrac{1}{4}\right)^{n-1}\right\}$である．

\newpage

\end{document}
